\section{A Levi reduction for a formal Higgs field}\label{sec:lie}

The Hecke modification in the next section must preserve a given Higgs
field and produce a nonzero covector in the direction of the moving point.
We obtain both properties from a cocharacter in the center of a Levi
subgroup.  The following proposition is the precise Lie-theoretic input.
Its two characteristic arguments are kept separate.

For a cocharacter $\lambda:\mathbb G_m\to G$, write
$d\lambda\in\mathfrak g$ for the image of $1\in\Lie(\mathbb G_m)=k$.
We distinguish the group center $Z(M)$ from the Lie-algebra center
$\mathfrak z(\mathfrak m)$, where $\mathfrak m=\Lie(M)$.
Write $\chi:\mathfrak g\to\mathfrak g/\!/G$ for the adjoint quotient.

\begin{proposition}\label{lie:input}
Let $k$ and $G$ satisfy Hypothesis~\ref{hyp:common} and either
Hypothesis~\ref{hyp:A} or Hypothesis~\ref{hyp:B}, and let $\kappa$ be
the invariant symmetric form in those hypotheses.
Suppose that $a(t)\in\mathfrak g[[t]]$ satisfies
$\chi(a(0))\ne\chi(0)$.  After a change of frame by an element of
$G(k[[t]])$, there are a Levi subgroup $M\subset G$ and a cocharacter
$\lambda:\mathbb G_m\to Z(M)$ such that
\begin{equation}\label{eq:lie-output}
 \begin{gathered}
 a(t)\in\mathfrak m[[t]],\qquad \kappa(a(0),d\lambda)\ne0,\\
 \operatorname{ad}_{a(0)}:\mathfrak g/\mathfrak m\longrightarrow
                   \mathfrak g/\mathfrak m\ \text{is invertible}.
 \end{gathered}
\end{equation}
Here $\operatorname{ad}_x(y)=[x,y]$.  The subgroup $M$ may equal $G$, and the
invertibility assertion includes the zero vector space.
\end{proposition}

The construction of a central modification follows the method of
\cite[\S\S20.4--20.5]{AGKRRV}.  We give the characteristic and lattice
arguments explicitly, since a vector in $\mathfrak z(\mathfrak m)$
alone does not specify a cocharacter of $Z(M)$.

\subsection{Root pairings and Jordan parts}

Fix a maximal torus $T\subset G$.  Put
$\Lambda=X_*(T)$, $\Lambda^\vee=X^*(T)$ and
$\mathfrak t=\Lie(T)=\Lambda\otimes_{\mathbb Z}k$.
For a root $\alpha$, its differential is the linear form
$d\alpha:\mathfrak t\to k$, whereas $d\alpha^\vee$ is the vector
obtained by differentiating its coroot.

\begin{lemma}\label{lie:root-pairing}
The form $\kappa$ is nondegenerate on $\mathfrak t$.  For every root
$\alpha$, there is a scalar $c_\alpha\in k^\times$ such that
\begin{equation}\label{eq:root-coroot-pairing}
 \kappa(d\alpha^\vee,h)=c_\alpha\,d\alpha(h)
 \qquad(h\in\mathfrak t).
\end{equation}
In particular, neither $d\alpha$ nor $d\alpha^\vee$ vanishes.
Every element of $\mathfrak g$ is conjugate into the Lie algebra of a
Borel subgroup containing $T$.  Every semisimple element is conjugate
into $\mathfrak t$.
\end{lemma}

\begin{proof}
The $T$-weight decomposition is
$\mathfrak g=\mathfrak t\oplus\bigoplus_\alpha\mathfrak g_\alpha$.
Invariance under $T$ makes $\mathfrak t$ orthogonal to the root spaces
and makes $\mathfrak g_\alpha$ pair only with
$\mathfrak g_{-\alpha}$.  Nondegeneracy of $\kappa$ therefore gives
nondegeneracy on $\mathfrak t$ and on each opposite-root pairing.
Choose compatible root vectors $e_\alpha,e_{-\alpha}$ with
$[e_\alpha,e_{-\alpha}]=d\alpha^\vee$.  This identity follows by
differentiating the rank-one root homomorphism from $\mathrm{SL}_2$;
it is valid for every isogeny type; see the root-subgroup construction
in \cite[\S\S8.1.1--8.1.4]{Springer} and the bracket formula in
\cite[Corollary~5.1.12, Equation~(5.1.4)]{Conrad}. Invariance of the form gives
\eqref{eq:root-coroot-pairing} with
$c_\alpha=\kappa(e_\alpha,e_{-\alpha})\ne0$.
Thus one of $d\alpha,d\alpha^\vee$ vanishes if and only if the other
does.  Their simultaneous vanishing would put
$\alpha\in p\Lambda^\vee$ and $\alpha^\vee\in p\Lambda$,
contradicting the integral identity
$\langle\alpha,\alpha^\vee\rangle=2$, since $p^2$ does not divide~$2$.
These arguments use torus characters, rather than just their
differentials, and also apply in characteristic two.

Let $B\supset T$ be a Borel subgroup with Lie algebra $\mathfrak b$.
The incidence morphism
\[
 G\times^B\mathfrak b\longrightarrow\mathfrak g,
 \qquad [g,x]\longmapsto\Ad(g)x,
\]
is proper: its source is closed in $(G/B)\times\mathfrak g$.
Choose $c\in\mathfrak t$ outside the root differential hyperplanes.
At $[1,c]$, the negative-root tangent directions map isomorphically
onto the negative root spaces, and $\mathfrak b$ supplies the remaining
directions.  The morphism is consequently dominant, hence surjective.

For a semisimple element in $\mathfrak b$, write it as $c+u$, with
$c\in\mathfrak t$ and $u$ in the positive-root Lie algebra.
Conjugate by positive root groups in increasing root height.  If
$d\alpha(c)\ne0$, conjugation by the $\alpha$-root group changes the
$\alpha$ coefficient by $-z\,d\alpha(c)$ and changes the other
positive-root terms only in greater height.  It therefore eliminates
that coefficient.  After all such steps the result is $c+u_0$ with
$[c,u_0]=0$.  In a faithful representation, $c$ is semisimple and
$u_0$ is nilpotent.  Since their sum is semisimple, the matrix Jordan
decomposition forces $u_0=0$.
\end{proof}

Choose a closed faithful representation $G\hookrightarrow\mathrm{GL}(V)$.
Its differential identifies $\mathfrak g$ with a restricted Lie
subalgebra of $\End(V)$: the operation $x\mapsto x^{[p]}$ is matrix
$p$th power.  This compatibility can also be defined intrinsically
by the $p$th iterate of a left-invariant derivation.  We use the
ordinary matrix Jordan decomposition, and below prove that its two
parts belong to $\mathfrak g$ in each regime; compare
\cite[\S4.4, Theorem~4.4.20]{Springer}.

\begin{lemma}\label{lie:constant-parts}
Suppose that $a_0=s+n\in\mathfrak g$, where $s,n\in\mathfrak g$,
$s$ is semisimple, $n$ is nilpotent in a faithful representation,
and $[s,n]=0$.  Suppose also that $\chi(a_0)\ne\chi(0)$.
After conjugation, $s$ is a nonzero element of $\mathfrak t$ and
$M=Z_G(s)$ is a Levi subgroup.  One has $a_0\in\mathfrak m$ and
$\operatorname{ad}_{a_0}$ is invertible on $\mathfrak g/\mathfrak m$.
For every $\lambda:\mathbb G_m\to Z(M)$,
\begin{equation}\label{eq:nilpotent-orthogonality}
 \kappa(a_0,d\lambda)=\kappa(s,d\lambda).
\end{equation}
\end{lemma}

\begin{proof}
Lemma~\ref{lie:root-pairing} conjugates $s$ into $\mathfrak t$.
If $s=0$, apply the Borel assertion of that lemma to $n$.
The projection $\mathfrak b\to\mathfrak t$ preserves restricted
$p$-operations.  Since $n$ is killed by an iterated $p$-operation
and that operation is injective on $\mathfrak t$, its toral
projection is zero.  A strictly dominant cocharacter contracts the
positive-root Lie algebra to zero.  Hence every invariant polynomial
has the same value on $n$ and on zero, a contradiction.  Thus $s\ne0$.

By Hypothesis~\ref{hyp:common}, $M=Z_G(s)$ is a Levi subgroup.
Its Lie algebra is the kernel of $\operatorname{ad}_s$, and it contains $a_0$.
On $\mathfrak g/\mathfrak m$, the operator $\operatorname{ad}_s$ is diagonalizable
with nonzero eigenvalues.  The commuting operator $\operatorname{ad}_n$ is
nilpotent, since
$\operatorname{ad}_n^{p^e}=\operatorname{ad}_{n^{[p^e]}}=0$ for large $e$.
On each eigenspace of $\operatorname{ad}_s$, the sum
$\operatorname{ad}_{a_0}=\operatorname{ad}_s+\operatorname{ad}_n$ is therefore invertible.

To prove \eqref{eq:nilpotent-orthogonality}, conjugate $n$ inside $M$
into the Lie algebra of a Borel subgroup of $M$ containing $T$.
The incidence argument of Lemma~\ref{lie:root-pairing} applies to
$M$: its roots have nonzero differentials, and $\kappa$ is
nondegenerate on $\mathfrak m$ by the same weight decomposition.
Restricted nilpotence again makes the toral projection of $n$ zero.
Thus the conjugated $n$ lies in the positive-root Lie algebra of $M$,
which is orthogonal to $\mathfrak t$.  Since $d\lambda\in\mathfrak t$
is fixed by $M$, invariance of $\kappa$ gives
$\kappa(n,d\lambda)=0$ before conjugation as well.
\end{proof}

It remains to choose the central cocharacter with nonzero pairing.
In the finite-field regime we prove an integral lattice statement;
in the arbitrary-field regime we use the separately assumed
invertibility of the Weyl-group order.

\subsection{The finite-field regime}

The lattice statement below uses only the common invariant-form and
Levi-center assumptions. We will then combine it with the finite-field
construction of the Jordan parts.

\begin{lemma}\label{lie:lattice}
Assume the invariant-form and Levi-center conditions of
Hypothesis~\ref{hyp:common}.  Every Levi Cartan determinant is prime
to $p$.  For every Levi subgroup $M$ containing $T$,
\begin{equation}\label{eq:central-span}
 \operatorname{span}_k\{d\lambda:
           \lambda\in X_*(Z(M))\}=\mathfrak z(\mathfrak m).
\end{equation}
Here $X_*(Z(M))$ denotes actual cocharacters of the group center,
viewed in $X_*(T)$.
\end{lemma}

\begin{proof}
After choosing a positive system, let
$I=\{\alpha_1,\ldots,\alpha_r\}$ be the simple roots of $M$.
Set $J_I=\operatorname{span}_k\{d\alpha_i^\vee\}\subset\mathfrak t$.
Because every root differential is nonzero, a central element of
$\mathfrak m$ has no root-space component.  Equation
\eqref{eq:root-coroot-pairing} therefore identifies
\[
 \mathfrak z(\mathfrak m)
   =\bigcap_i\ker(d\alpha_i)=J_I^\perp.
\]
The form is nondegenerate on this center by hypothesis, and on
$\mathfrak t$ by Lemma~\ref{lie:root-pairing}; it follows that it is
nondegenerate on $J_I$.

Define integral maps
\[
 U:\Lambda\longrightarrow\mathbb Z^r,\quad
       y\longmapsto(\langle\alpha_i,y\rangle)_i,
 \qquad
 V:\mathbb Z^r\longrightarrow\Lambda,\quad
       e_j\longmapsto\alpha_j^\vee.
\]
Their composite is the Cartan matrix
$C_I=UV=(\langle\alpha_i,\alpha_j^\vee\rangle)_{i,j}$.
A subscript $k$ will mean reduction to $k$.
Equation~\eqref{eq:root-coroot-pairing} gives
$V_k^*\kappa=D U_k$, where
$D=\operatorname{diag}(c_{\alpha_1},\ldots,c_{\alpha_r})$ is invertible.
The radical of $V_k^*\kappa V_k$ is exactly $\ker V_k$, by
nondegeneracy on $J_I=\im V_k$.  Consequently
\begin{equation}\label{eq:cartan-ranks}
 \operatorname{rank}(C_{I,k})
  =\operatorname{rank}(V_k)=\operatorname{rank}(U_k).
\end{equation}

If some Levi Cartan determinant were divisible by $p$, the ordinary
determinant list would give a further Levi of type $A_{p-1}$.
Indeed, the Cartan matrices of the classified Dynkin diagrams
\cite[Theorem~C.56 and \S24(c)]{Milne} have determinants
\[
 \begin{array}{c|ccccccccc}
 \text{type}&A_r&B_r&C_r&D_r&E_6&E_7&E_8&F_4&G_2\\ \hline
 \det&r+1&2&2&4&3&2&1&1&1.
 \end{array}
\]
In type $A_r$, take a consecutive string of $p-1$ vertices.
In the determinant-two and determinant-four cases take one vertex,
since $p=2$.  In type $E_6$ take two adjacent vertices, since $p=3$.
Such a simple-root subdiagram defines a Levi of $G$, so the preceding
rank calculation applies to it.

For this $A_{p-1}$ Levi, $r=p-1$ and $\det C_I=p$.
Choose dual bases of the full character and cocharacter lattices,
of rank $N_T$.  The matrices of $U,V$ have sizes
$r\times N_T$ and $N_T\times r$.  By
\eqref{eq:cartan-ranks}, both reductions have rank less than $r$,
so every maximal minor of each matrix is divisible by $p$.
The integral Cauchy--Binet identity gives
\[
 \det C_I=
 \sum_{\substack{J\subset\{1,\ldots,N_T\}\\ |J|=r}}
       \det(U_{:,J})\det(V_{J,:}),
\]
which is divisible by $p^2$, contrary to $\det C_I=p$.
Using the full lattices here retains all central torus directions
and makes no assumption on the isogeny type.

For an arbitrary Levi, put $d_I=\det C_I$ and define an integral
endomorphism of $\Lambda$ by
\[
 \Pi_I=d_I\id_\Lambda-V\operatorname{adj}(C_I)U.
\]
It satisfies $U\Pi_I=0$.  The lattice $\ker U$ consists precisely of
cocharacters into $Z(M)$: their images commute with $T$ and with
every root group of $M$.  For $z\in\ker U_k$,
$\Pi_{I,k}(z)=d_Iz$.  Since $d_I$ is invertible in $k$, the
differentials of elements of $\ker U$ span
$\ker U_k=\mathfrak z(\mathfrak m)$, proving
\eqref{eq:central-span}.  For $r=0$, take $d_I=1$ and
$\Pi_I=\id_\Lambda$; the same proof covers tori.
\end{proof}

We now assume Hypothesis~\ref{hyp:A}, so
$k=\overline{\mathbb F}_q$, and construct the constant-term data in
Proposition~\ref{lie:input}.  In a faithful representation, write
$a_0=a(0)=s+n$ for its matrix Jordan decomposition.  The eigenvalues
of $s$ lie in a finite field $\mathbb F_{p^h}$.  Choose $e$ divisible
by $h$ and large enough to have $n^{p^e}=0$.  Then
\[
 a_0^{[p^e]}=a_0^{p^e}=s^{p^e}+n^{p^e}=s.
\]
Thus $s,n\in\mathfrak g$.  Lemma~\ref{lie:constant-parts} supplies
a nonzero $s\in\mathfrak t$, the Levi $M=Z_G(s)$ and invertibility
of $\operatorname{ad}_{a_0}$ on $\mathfrak g/\mathfrak m$.
The element $s$ lies in $\mathfrak z(\mathfrak m)$, where $\kappa$
is nondegenerate.  By Lemma~\ref{lie:lattice}, some actual
$\lambda\in X_*(Z(M))$ satisfies $\kappa(s,d\lambda)\ne0$.
Equation~\eqref{eq:nilpotent-orthogonality} gives the required
$\kappa(a_0,d\lambda)\ne0$.  This construction has used no
Weyl-order assumption.

\subsection{The arbitrary algebraically closed regime}

We now assume Hypothesis~\ref{hyp:B}.  The eigenvalues of a matrix
over $k$ need not be fixed by any positive power of Frobenius.
Instead, a finite additive-polynomial interpolation gives its
Jordan parts inside $\mathfrak g$.

Write $a_0=s+n$ in a faithful representation and choose $e$ with
$n^{p^e}=0$, so $a_0^{p^e}=s^{p^e}$.  Let $V_0\subset k$ be the
finite-dimensional $\mathbb F_p$-span of the eigenvalues of $s$,
and let $V_e=\{v^{p^e}:v\in V_0\}$.  The inverse of Frobenius
$V_e\to V_0$ is $\mathbb F_p$-linear.  If $V_0=0$, then $s=0$ and
$a_0=n\in\mathfrak g$, contrary to
Lemma~\ref{lie:constant-parts} and $\chi(a_0)\ne\chi(0)$.
We may therefore assume that $r=\dim_{\mathbb F_p}V_e\ge1$.
If
$\beta_1,\ldots,\beta_r$ is an $\mathbb F_p$-basis of $V_e$, the
Moore matrix
\[
 (\beta_i^{p^j})_{1\le i\le r,\ 0\le j<r}
\]
is invertible.  Otherwise a nonzero polynomial
$\sum_{j=0}^{r-1}c_j z^{p^j}$ of degree at most $p^{r-1}$ would
vanish on all $p^r$ elements of $V_e$, a contradiction.
There is therefore an additive polynomial
$Q(z)=\sum_{j=0}^{r-1}c_jz^{p^j}$ agreeing with inverse Frobenius
on $V_e$.  Applying it to the diagonalizable matrix $s^{p^e}$ gives
\[
 s=Q(a_0^{p^e})=\sum_{j=0}^{r-1}c_j a_0^{[p^{e+j}]}\in\mathfrak g.
\]
Thus $n\in\mathfrak g$ as well.  Lemma~\ref{lie:constant-parts} gives
$s\in\mathfrak t\setminus\{0\}$, $M=Z_G(s)$ and the quotient
invertibility.

Let $W_M=N_M(T)/T$ be the Weyl group of $M$ and set $h_M=|W_M|$.
Since $W_M\subset W_G$, the integer $h_M$ is invertible in $k$.
The operator
\[
 e_M=\frac1{h_M}\sum_{w\in W_M}w:\mathfrak t\longrightarrow\mathfrak t
\]
is a self-adjoint projection for $\kappa$.
Consequently $\kappa$ is nondegenerate on
$\mathfrak t^{W_M}=\im e_M$.  This subspace contains $s$, since
$M$ centralizes $s$.
In fact $\mathfrak t^{W_M}=\mathfrak z(\mathfrak m)$:
the differential reflection formula is
$s_\alpha(h)=h-d\alpha(h)d\alpha^\vee$, and
$d\alpha^\vee\ne0$ by Lemma~\ref{lie:root-pairing}.
Thus, under the Weyl-order assumption, nondegeneracy on every Levi
Lie center also follows from the invariant form itself.

To obtain integral cocharacters, choose a basis
$\mu_1,\ldots,\mu_{N_T}$ of $\Lambda$ and form the integral sums
\[
 \nu_i=\sum_{w\in W_M}w\mu_i\in\Lambda^{W_M}.
\]
Their differentials span $\mathfrak t^{W_M}$ because
$d\nu_i=h_Me_M(d\mu_i)$.  The integral reflection formula
\[
 s_\alpha(\lambda)=\lambda-
             \langle\alpha,\lambda\rangle\alpha^\vee
\]
and torsion-freeness of $\Lambda$ show that
$\Lambda^{W_M}=X_*(Z(M))$: an invariant cocharacter has zero
integral pairing with every root of $M$, and conversely.
No division of a cocharacter by $h_M$ has been made.
Nondegeneracy on $\mathfrak t^{W_M}$ now gives an integral central
cocharacter $\lambda$ with $\kappa(s,d\lambda)\ne0$.
By \eqref{eq:nilpotent-orthogonality},
$\kappa(a_0,d\lambda)\ne0$ as required.

This argument applies to arbitrary algebraically closed $k$.
Its use of $p\nmid|W_G|$ is specific to Hypothesis~\ref{hyp:B};
the finite-field argument above obtained its cocharacter from the
Levi-center form condition instead.

\subsection{A single formal reduction}

Both regimes have now provided the same constant-term data.
The remaining step moves the entire formal Higgs coefficient into
the chosen Levi.  It uses neither finite-field descent nor a
Weyl-group average.

\begin{lemma}\label{lie:formal}
Let $G$ be a connected reductive group over an algebraically closed
field $k$, let $M\subset G$ be a Levi subgroup, and let
$a(t)\in\mathfrak g[[t]]$ satisfy $a_0=a(0)\in\mathfrak m$.
If $\operatorname{ad}_{a_0}$ is invertible on $\mathfrak g/\mathfrak m$, then
there exists $g(t)\in G(k[[t]])$ with $g(0)=1$ such that
$\Ad(g(t))a(t)\in\mathfrak m[[t]]$.
\end{lemma}

\begin{proof}
Suppose that the current arc belongs to $\mathfrak m[[t]]$
modulo $t^r$, with $r\ge1$, and let $c_r\in\mathfrak g/\mathfrak m$
be the class of its coefficient of $t^r$.  Choose $D\in\mathfrak g$
such that $c_r+[D,a_0]=0$ in the quotient.
Smoothness of $G$ gives an element $g_r(t)\in G(k[[t]])$ with
\[
 g_r(t)\equiv1+t^rD\pmod {t^{r+1}}.
\]
Here the displayed expression specifies a point in the tangent
kernel for the square-zero ideal $(t^r)/(t^{r+1})$; formal
smoothness lifts it successively to every higher order.
Conjugation by $g_r(t)$ changes the coefficient of $t^r$ by
$[D,a_0]$ and leaves all lower coefficients unchanged.
It therefore places the arc in $\mathfrak m[[t]]$ modulo $t^{r+1}$.

Iterate this construction.  Since $g_r(t)\equiv1\pmod {t^r}$,
the successive products converge in the $t$-adic topology.
The affine group $G$ takes this compatible system of points modulo
$t^j$ to a point $g(t)\in G(k[[t]])$ with $g(0)=1$.
Its conjugation sends the arc into $\mathfrak m[[t]]$.
No exponential series is used.
\end{proof}

\begin{proof}[Proof of Proposition~\ref{lie:input}]
Under Hypothesis~\ref{hyp:A}, use the restricted-power construction
and Lemma~\ref{lie:lattice}.  Under Hypothesis~\ref{hyp:B}, use the
additive-polynomial construction and integral Weyl-group averaging.
In each case Lemma~\ref{lie:constant-parts} supplies the Levi and
quotient invertibility, with
$\kappa(a_0,d\lambda)\ne0$.  Apply Lemma~\ref{lie:formal} after the
initial constant conjugation.  The formal change of frame fixes
$a_0$, so all three assertions in \eqref{eq:lie-output} hold.
\end{proof}

\begin{remark}\label{lie:zero-orbit}
Centrality of $\lambda$ in $M$ makes
$\Ad(\lambda(t))a(t)=a(t)$.  Moreover, every nonzero
$\lambda$-weight space of $\mathfrak g$ injects into
$\mathfrak g/\mathfrak m$, so the adjoint action is invertible on
the truncated nonzero-weight modules used below.
We do not need $Z_G(\lambda)=M$.  If $\lambda$ is central in $G$,
the Hecke orbit has dimension zero; the nonzero pairing in
\eqref{eq:lie-output} still supplies the moving-point covector.
In particular, tori are included in every allowed characteristic.
\end{remark}
